\section{Sim-to-Real 与验证实验}

\subsection{仿真校准与现实验证}

\begin{figure}[H]
\centering
\includegraphics[width=0.8\textwidth]{slide-02-11.jpg}
\caption{Slide 02-11 说明在 sim-to-real pipeline 中如何通过 domain randomization、real world replicas 来校准。}
\end{figure}

为了让 policy 从仿真迁移到真实环境，常用的策略包括 domain randomization、dynamics randomization、fidelity tuning。每次迁移都要记录 randomization 参数与 reality gap 监控指标。

\begin{table}[H]
\centering
\begin{tabular}{p{0.3\textwidth} p{0.6\textwidth}}
\toprule
验证点 & 说明 \\
\midrule
Domain gap & 动力学/感知差异的 distribution drift score \\
Simulator fidelity & 视觉/控制 fidelity 与 real world log 的 L2 误差 \\
Transfer success & real world rollout 的 task success rate \\
\bottomrule
\end{tabular}
\caption{Sim-to-Real 验证矩阵}
\end{table}

\begin{knowledgebox}{记录 transfer meta-data}
用 metadata 记录每次 transfer 的 simulator seed、randomization range、hardware configuration，有助于复现某次 jump 的成功或 failure。
\end{knowledgebox}

\subsection{Cross-environment evaluation}

不同环境（lab、field、cloud）会有不同 constraints。需要规定 evaluation order，例如：先在 lab environment run benchmark，再到 field environment 检查 sensor noise 和 drift，再最后在 cloud pipeline 里收集 logs。

\begin{warningbox}{不要在 field rollouts 上直接用 baseline}
Field environment 一旦 drift，就非常 expensive。要在 lab 中先验证 metric，然后 incremental expansion 到 field，最后在 cloud 把 logs 收集到 evidence board 中。
\end{warningbox}

\subsection{本章小结}

Sim-to-Real pipeline 要靠 structured validation matrix 与 metadata，才能在不同环境之间安全迁移 policy；在每次 field rollout 之前必须验证 domain gap。

